\documentclass[10pt,a4paper]{article}
\usepackage[margin=18mm]{geometry}
\usepackage{amsmath,amssymb,booktabs,array}
\usepackage{cite}
\renewcommand{\citepunct}{], [}
\renewcommand{\citedash}{]--[}
\input{glyphtounicode}
\pdfgentounicode=1
\usepackage[hidelinks]{hyperref}
\hypersetup{pdftitle={An urban signal-control task: model and sensitivity},
 pdfauthor={Urban traffic methods},pdfsubject={Compact methodology}}
\setlength{\parindent}{0pt}
\setlength{\parskip}{3pt}
\setlength{\abovedisplayskip}{4pt}
\setlength{\belowdisplayskip}{4pt}
\setlength{\abovedisplayshortskip}{3pt}
\setlength{\belowdisplayshortskip}{3pt}
\setlength{\emergencystretch}{2em}
\widowpenalty=10000
\clubpenalty=10000
\newcommand{\pos}[1]{[#1]_+}
\newcommand{\head}[1]{\par\vspace{4pt}\noindent\textbf{#1}\par\nobreak\vspace{1pt}}
\newenvironment{points}{\begin{list}{$\bullet$}{%
 \setlength{\leftmargin}{1.4em}\setlength{\labelwidth}{0.7em}%
 \setlength{\labelsep}{0.4em}\setlength{\itemsep}{2pt}%
 \setlength{\parsep}{0pt}\setlength{\topsep}{3pt}}}{\end{list}}
\usepackage{xcolor}
\definecolor{revisionred}{HTML}{D13F43}
\newcommand{\rev}[1]{{\color{revisionred}#1}}
\begin{document}
\begin{center}
{\large\bfseries An urban signal-control task: model and sensitivity}
\end{center}
\vspace{-5pt}

\head{1. Research questions}
A signal can offer more left-turn service while a turner remains unable
to reach the bay. The relevant obstruction depends on which traffic is
in front, which lane it occupies, and whether another lane remains usable.
The task asks how much of that ordered information is needed for useful
signal-control sensitivities:
\begin{points}
\item \textbf{Information:} Can retained lane order explain reachable
service beyond sound aggregate models using movement counts, storage
and their own inference from traffic histories?
\item \textbf{Decision:} Does that additional information change which
effective green allocation reduces total time for the complete demand?
\item \textbf{Gradient:} Does the predicted response to a small service
reallocation agree with actual changes and help choose a better control?
\end{points}
These are separate claims: a more detailed state can improve a forward
prediction without necessarily providing a more useful control direction.

\head{2. Task and global assumptions}
\textbf{Control and population.} Consider one isolated, initially symmetric
four-leg junction with a finite set of releases and fixed routes.
One constant vector $g\geq0$, $\sum_{p=1}^4g_p=1$, allocates effective
service to NS through/right, NS left, EW through/right and EW left.
Green fractions represent average service, rather than the active phase
at a particular instant. Traffic is loaded from empty; any preparation
period is the same for every candidate control. After the control window,
a common fixed allocation clears all released demand.

\textbf{Representative approach.} There are two lanes before the bay
opening and three at the stop line. Aggregate upstream road $U$ feeds
parallel local lanes $F_1$ (left/through) and $F_2$ (through/right).
At the opening, $F_1$ branches into left-only bay $B$ and through-only
continuation $T_1$; $F_2$ continues into shared through/right lane $T_2$.

ev{The connectivity below defines this task's representative geometry:}
\[
 U\longrightarrow\{F_1,F_2\},\qquad
 F_1\longrightarrow\{B,T_1\},\qquad F_2\longrightarrow T_2.
\]
The bay opening lies inside the local zone. Its cross-section also
separates $F_2$ from $T_2$, so the two feeders cover the same upstream
length. In the symmetric core $B,T_1,T_2$ have the same remaining length.
Every section owns distinct road space and storage.

Left traffic enters $F_1$ and must use $B$; right traffic uses $F_2,T_2$.
Through traffic can choose $F_1$ or $F_2$ when first admitted from $U$.
Once in a feeder it stays in that lane: $F_1$ through traffic continues
to $T_1$, and $F_2$ through traffic continues to $T_2$.
The admission solver compares current feeder loads and feasible entry;
it does not move resident traffic between lanes or predict a future
signal plan. A blocked $F_1$ therefore need not stop traffic already
using $F_2$. An occupied $F_1$ head still blocks traffic behind it.
The smaller single-feeder layout remains a mechanism diagnostic;
compared models always receive the same declared geometry and permissions.

\textbf{Modeling assumptions.}
\begin{points}
\item \textbf{MA1: continuous quantity, fixed routes.} Traffic mass is
continuous vehicle-equivalent quantity. Routes, movements, lane permissions
and topology are fixed metadata. Each realized lane retains an ordered
record. Its cumulative-mass coordinate measures order, not physical metres.
\item \textbf{MA2: average signal service.} Green fractions change movement
service and downstream sharing. Road capacities, travel times, backward
wave speeds, storage and saturation parameters are fixed. The model does
not resolve red/green switching within a cycle.
\item \textbf{MA3: physical access.} An unassigned input list determines
which packet the allocator examines first. A packet already occupying a
lane also has a physical position: traffic behind a blocked head cannot
pass. Independent movement requires an independently reachable lane.
\item \textbf{MA4: entry establishes lane occupancy.} Each packet is
allocated using current receiving-lane loads and known physical limits.
Later packets see earlier actual admissions. Unadmitted mass remains
where it is, with no binding future-lane target and with priority over
later packets of its movement.
\item \textbf{MA5: causal timesteps.} Sending and receiving limits use
histories available at the start of the update. New inflow cannot leave
the same section in that update, and current departures do not instantly
create entrance room. Boundary histories are interpolated linearly.
\item \textbf{MA6: complete demand accounting.} Releases, routes and
preparation are fixed across controls. Source waiting, moving traffic and
queues all contribute until every released trip clears. Queue differences
continue to matter after the control window ends.
\end{points}

\par\begin{samepage}
\textbf{Travel and storage.} For each section $e$, capacity $C_e$ is a
mass rate and $K_e$ is maximum stored mass. Free-flow travel time
$\tau^f_e$ is the time to cross the section downstream. Backward-wave
time $\tau^b_e$ is the time for a congestion or space-release effect to
cross upstream. The storage identity below is
\rev{derived from the triangular FD in the project's LTM transcription~\cite{m0}};
the timestep bound is \rev{the numerical domain adopted in that transcription}:
\begin{equation}
 K_e=C_e(\tau^f_e+\tau^b_e),\qquad
 0<\Delta t\leq\min_e\{\tau^f_e,\tau^b_e\}.
 \label{eq:compact-domain}
\end{equation}
For example, downstream departures release space, but the entrance can
use that space only after the backward wave reaches it. Numerical
lengths, demands and clearance horizons must satisfy this domain;
they are not calibrated by the equations alone.
\par\end{samepage}

\head{3. Traffic model}
The model follows a traffic journey. First determine which resident
mass has travelled far enough and where new mass can fit. Then assign
accessible traffic to lanes. Stop-line heads use their movement's service
and compete for downstream room. Finally update each physical section
and retain the histories needed for the next timestep.
Cumulative propagation and ordered route bookkeeping provide the
backbone~\cite{m0,m1}; the lane allocation and access rules below specify
the local urban interaction.

\textbf{3.1. Propagation and eligible arrivals.}
A section can send traffic only after it has travelled far enough, and
can receive traffic only when room is available at its entrance.
At $t_k=k\Delta t$, let $N^+_{e,k}$ count all earlier entries and
$N^-_{e,k}$ all earlier departures. Their difference
$n_{e,k}=N^+_{e,k}-N^-_{e,k}$ is resident mass.
The functions ${\cal H}^{\pm}_e$ read those cumulative histories between
stored boundaries by linear interpolation:
${\cal H}^{\pm}_e((k+\theta)\Delta t)
=(1-\theta)N^\pm_{e,k}+\theta N^\pm_{e,k+1}$.
Histories are zero before empty loading.
\rev{This interpolation reuses the project's discrete history reconstruction~\cite{m0,m1}.}

The sending calculation first counts entries early enough to reach the
exit by $t_{k+1}$ under free flow, then subtracts mass that has already
left. This gives eligible resident mass $A$. Only capacity $C_e\Delta t$
can cross during the update, so sending $S$ is the smaller amount.
Receiving uses the corresponding space account: start with storage
$K_e$, add departures whose space-release effect can reach the entrance
by $t_{k+1}$, and subtract all previous entries.
\rev{The limits below reuse the project's discrete transcription of Yperman's LTM~\cite{m0}.}
Applying them to the internal urban sections changes the section graph,
not the propagation law.
\begin{align}
 A_{e,k}&=\pos{{\cal H}^+_e(t_{k+1}-\tau^f_e)-N^-_{e,k}},
 &S_{e,k}&=\min(C_e\Delta t,A_{e,k}),\label{eq:compact-sending}\\
 R_{e,k}&=\min\{C_e\Delta t,\,
 \pos{K_e+{\cal H}^-_e(t_{k+1}-\tau^b_e)-N^+_{e,k}}\}.
 \label{eq:compact-receiving}
\end{align}
Here $[z]_+=\max(z,0)$. In the first line, the delayed entrance history
sets travel-time eligibility; in the second, the delayed exit history
sets entrance-accessible room. Both $S$ and $R$ are mass allowances for
this update, not instantaneous flow rates or guaranteed actual transfers.
Signals, access and neighbouring sections can reduce realized flow.
The timestep condition makes both delayed queries no later than $t_k$.

\textbf{3.2. Lane allocation and physical access.}

Lane-equilibrium models explain lane use through the resistance of
permitted alternatives; Gong, Bliemer and Raadsen formulate equilibrium
allocation in a lane-based node model~\cite{gong}. Here the resistance is
current queue mass. For each arriving homogeneous packet, allocate as much
as can physically enter, then balance the receiving loads. This is a
queue-based entry adaptation of that idea.
\rev{The capped entry formula below is derived here from constrained quadratic load balancing; Gong et al. supply its equilibrium motivation, not this entry equation.} Existing residents stay in place. A packet carries a fixed
route reference, current movement and continuous mass/timing data~\cite{m1}.
Section 3.5 explains the linked ledger and grid representations.
\par

For this packet, $q$ is its eligible mass and $s$ the sending owner's
remaining transfer allowance. For each legal, reachable receiving lane,
$u_\ell$ is its remaining admission allowance and $h_\ell$ its current
load estimate: start-of-step resident mass plus earlier actual admissions.
Thus $V$ answers ``how much can enter in total?'' It is limited by the
packet, the source and the combined receiving room. The lane solver then
chooses a common filling level $\lambda$ and additions $x_\ell$:
\begin{equation}
 V=\min(q,s,\sum_\ell u_\ell),\qquad
 x_\ell=\min\{u_\ell,[\lambda-h_\ell]_+\},\quad
 \sum_\ell x_\ell=V .
 \label{eq:compact-allocation}
\end{equation}
The unknown $\lambda$ is a queue-load level, measured in vehicle-equivalent
mass. The gap $\lambda-h_\ell$ is how much that lane would need to reach
the level. Add nothing to a lane already above it, and never exceed its
individual bound $u_\ell$. Choose the level so that all additions sum to
$V$. This is the specified ideal lane-balancing rule, not a stop-line
service cap. With initial loads $0.2,0$, admitted mass $V=0.6$ and
nonbinding lane bounds, $\lambda=0.4$ gives additions $0.2,0.4$.
A capped lane can remain below the common level. The additions are unique;
$\lambda$ itself can be nonunique when every usable lane is capped.

In the representative geometry, through packets have two alternatives
at $U\to\{F_1,F_2\}$. Left and right packets have one legal feeder each,
and movement fixes the branch at the bay opening. The same equation
therefore covers both responsive choice and forced entry. Lane-entry
room is separate from stop-line service: zero green does not immediately
prevent a vehicle from joining a lane that it can reach and that has room.

Only admitted amounts establish new lane occupancy. Update the receiving
loads by $h\leftarrow h+x$ and subtract those transfers from the remaining
sending and receiving allowances. The remainder $q-V$ stays at its current
location without a binding future-lane target. Inspect packets in recorded
order; if one movement remains partly unserved, its later packets wait,
while other movements with physical access can still be examined.

\par\begin{samepage}
The sending allowance limits total mass actually admitted, not how many
ready packets may be inspected. For example, an unassigned pool may list
one blocked left-turn unit before a through unit with an independently
reachable lane. With sending allowance one, inspecting only the first
unit would miss the feasible through admission. Examining the through
packet uses the same one-unit limit without moving it past a physical
blocker. Unserved mass keeps its original entry/travel history and its
relative position among the remaining entries; waiting does not restart
its travel delay.
\par\end{samepage}

Physical lane order remains binding. A residual at the front of $F_1$
prevents following traffic from reaching either $B$ or $T_1$, even if
the allocator examines its record. Traffic already in $F_2$ has its own
path to $T_2$. No traffic can jump from an occupied feeder into the
parallel lane or enter $B,T_1$ directly from $U$.


\textbf{3.3. Stop-line service and downstream sharing.}
A green fraction represents average service opportunity, following the
continuum-signal idea of Han et al.~\cite{han}.
\rev{The law below adapts that average-service interpretation to this task's four groups and physical-cap bound.}
Saturation rate $s_{\ell m}$ times fraction $g_{p(m)}$ supplies the
uncapped rate of a ready lane head.
\par
\begin{equation}
 g_p\geq0,\qquad \sum_{p=1}^4g_p=1,\qquad
 \mu_{\ell m}(g)=\min\{C_\ell,s_{\ell m}g_{p(m)}\}.
 \label{eq:compact-action}
\end{equation}

The head's rate cap is $c_\ell=\mu_{\ell m}$, in vehicle-equivalents/s.
Downstream receiving can reduce its usable rate. Write
$\rho_j=R_{j,k}/\Delta t$ for output $j$'s average receiving rate.
The demand/supply and conservation interpretation follows the first-order
node framework~\cite{node}.
\rev{That reference supplies feasibility and conservation requirements; the two sharing laws below are defined for this model.}
The priority rule gives
each competing inbound leg a green-weighted share, then divides that
share equally among its ready lanes, reusing portions that a capped lane
cannot use.

Let ${\cal L}_{ij}$ be the ready lanes from inbound leg $i$ whose current
head enters output $j$, with positive remaining sending. Their total rate
cap is $C_{ij}$; $w_{ij}$ is the relevant green weight. The usable total
output rate is $W_j$. The solver finds $\theta_j$, a common multiplier of
leg weights, so leg rates sum to $W_j$. It then finds $\eta_{ij}$, the
common uncapped lane rate within leg $i$ toward $j$:
\par
\begin{align}
 C_{ij}&=\sum_{\ell\in{\cal L}_{ij}}c_\ell,\quad
 w_{ij}=g_{p(i,j)},\quad W_j=\min(\rho_j,\sum_iC_{ij}),\nonumber\\
 r_{ij}&=\min(C_{ij},\theta_jw_{ij}),\quad
                    \sum_i r_{ij}=W_j,\label{eq:compact-leg-sharing}\\
 r_\ell&=\min(c_\ell,\eta_{ij}),\quad
                    \sum_{\ell\in{\cal L}_{ij}}r_\ell=r_{ij}.
 \label{eq:compact-lane-sharing}
\end{align}

Thus $r_{ij}$ is the leg's granted rate and $r_\ell$ the lane's actual
rate. The unknown $\eta_{ij}$ is measured in vehicle-equivalents/s; each
lane receives the smaller of it and its cap $c_\ell$. For a leg grant
$0.6$ shared by two unrestricted lanes, $\eta_{ij}=0.3$. If one lane is
capped at $0.1$, the other can use $0.5$, giving $\eta_{ij}=0.5$.
A leg is green-weighted once, rather than once for every lane it contains.
These priorities specify the present model; they are not a transcription
of the oriented-capacity priorities in the node reference.

The sharing equations give rates for the current heads. The simulation
needs amounts for the whole timestep. A head can clear before that time
is used up, exposing a different movement with a different rate and
possibly a different downstream competitor set.

For a complete example, take a shared left/through diagnostic lane with
one second available, a ready left block of mass $0.2$ followed by at
least $0.3$ ready through mass, and ample sending and downstream room.
At average left rate $0.4$, clearing the front block uses half a second.
At through rate $0.6$, the next block then discharges $0.3$ in the remaining
half. The one-step output is therefore $0.2$ left plus $0.3$ through.
The block order determines how the timestep's service time is used.

For several interacting lanes, the earliest head change can alter their
shared receiving rates. Let $\sigma$ be time already accounted for, starting
at zero; $q_\ell$ the current head's remaining mass; and
$S_\ell^{\rm rem}$ the lane's remaining sending allowance. The duration
$\delta$ for which the present rates continue is the smallest of remaining
timestep time and the time to clear a head or spend a sending allowance.
\par\begin{samepage}
\rev{Combining these task-defined stopping constraints derives the minimum below:}
\par
\begin{equation}
 \delta=\min\left\{\Delta t-\sigma,\
 \min_{\ell:r_\ell>0}
       \frac{\min(q_\ell,S_\ell^{\rm rem})}{r_\ell}\right\}.
 \label{eq:compact-head-time}
\end{equation}
\par\end{samepage}


Each lane contributes its rate times this duration. After a head clears,
its next already-resident eligible block uses its own rate and competes
with the other ready heads for the remaining time. A head with zero
usable rate leaves its followers waiting. For a chosen lane and movement,
let $a$ index these successive intervals and let $r_a$ be its granted
rate, zero when that movement is not the head.
\rev{Summing mass transferred at each rate over its duration derives the timestep amount:}
\begin{equation}
 y=\sum_a r_a\delta_a.
 \label{eq:compact-service-amount}
\end{equation}
The final interval covers the remaining timestep time, with zero rate
for an idle stream. The transferred amounts, rather than the intermediate
intervals, enter the section update below.
\par

\textbf{3.4. Mass conservation and travel histories.}
Every internal transfer must leave one physical section and enter the
next with the same mass and route identity. If $b_{e,k}$ enters and
$y_{e,k}$ leaves section $e$,
\rev{the inherited cumulative-count bookkeeping~\cite{m0,m1} gives the conservation update:}
\begin{equation}
 N^+_{e,k+1}=N^+_{e,k}+b_{e,k},\quad
 N^-_{e,k+1}=N^-_{e,k}+y_{e,k},\quad
 n_{e,k+1}=n_{e,k}+b_{e,k}-y_{e,k}.
 \label{eq:compact-update}
\end{equation}
The source first adds released mass $D_{o,k}$ to its waiting mass $Q_{o,k}$;
admission then moves it onto the road. Sink exit removes it from the
road population. The bounds $b\leq R\leq K-n$ and $y\leq S\leq n$
keep occupancy between zero and storage. Unadmitted mass is counted
only where it physically remains. New admissions are appended after
resident service has been selected.


A stored lane count alone cannot tell which mass has had time to reach
its exit. The route ledger therefore records when each admitted portion
entered its new section. This lets a later update identify the eligible
part of each ordered packet, and keeps admission timing consistent with
the cumulative propagation history.

\rev{Reuse the project's uniform boundary-history reconstruction~\cite{m1}:} spread total
entrance mass $b>0$ across the timestep of length $\Delta t$, beginning at
$t_k$, in actual admission order. Consider an admitted packet of mass
$q^{\rm in}$ with $z$ admitted mass before it. The two endpoints below
are both \emph{entry} times: when its first and last parts enter.
\rev{Their expressions follow from locating the packet's mass fractions within this uniformly admitted stream.}
\begin{equation}
 [t^{\rm in}_{\rm front},t^{\rm in}_{\rm tail}]
 =\left[t_k+\Delta t\,\frac z b,\
       t_k+\Delta t\,\frac{z+q^{\rm in}}b\right].
 \label{eq:compact-entry}
\end{equation}
The fraction $z/b$ locates its beginning within the admitted stream;
$(z+q^{\rm in})/b$ locates its end. With one unit entering during a
one-second update, $0.2$ before a $0.3$-unit packet assigns that packet
entry from $0.2$ to $0.5$ seconds after $t_k$.

The free-flow travel time $\tau^f_e$ is the time to cross section $e$
when unobstructed. Adding it to both endpoints gives when the packet's
first and last parts could reach the exit. If a later update cuts through
this readiness interval, only the corresponding part is eligible to send.
Waiting traffic retains that history, so a failed admission does not
restart its travel time. This is how packet identity and propagation
availability stay connected.
\par


\textbf{3.5. Link ledgers, lane ledgers and the mass grid.}
The model retains identity and quantity in two connected forms.
A long-link ledger is an ordered list of route-labelled continuous
packets. Once traffic enters the refined zone, each realized lane has
its own list. A record says which fixed route the packet follows, its
current place on that route, its mass and its readiness-time interval.
When part moves, split its mass continuously and leave the remainder
in order. Entering a turn bay still leaves the packet on the same incoming road;
its position in the global route advances only when it actually crosses
the junction onto the next road.

Tracking identity on an LTM backbone is also the central idea in the
individual-tracking extension of de Souza et al.~\cite{tracking}.
\rev{That work is a conceptual precedent for identity tracking, not the source of the continuous ledger or grid equations here.}
Their integer-flow construction is a useful contrast: here quantities
remain continuous and the route/order labels are retained as metadata.
The ledger construction and continuous split/timing rules are specific
to this model~\cite{m1}. A small service change can therefore change a
packet's transferred or remaining mass without waiting for a whole
vehicle to move.

The near-intersection grid supplies the per-lane continuous movement
state. It is cell-based, but its cells are intervals of cumulative mass
from the lane head, rather than fixed road-distance cells as in the
classical CTM~\cite{ctm}. LTM histories and geometry still determine
physical travel and storage. This mass-coordinate adaptation represents
local composition without adding another physical wave-propagation law.
\rev{The formula below is derived from packet/bin interval intersections in the project's mass coordinates; CTM supplies the cell-based inspiration.}

\par\begin{samepage}
Let packet $r$ in lane $\ell$ occupy mass-rank interval
$[z^-_{\ell r},z^+_{\ell r}]$, with movement $m_r$ obtained from its
route. Fixed grid bin $c$ spans
$[\zeta_{\ell c},\zeta_{\ell,c+1}]$. Bin edges respect the lane's
physical storage, including any partial final bin. Let $\omega_{\ell cr}$
be the mass of packet $r$ lying inside bin $c$ of lane $\ell$.
Then $M_{\ell cm}$ is movement $m$'s mass in that bin, while $n_\ell$
is total resident lane mass, including both moving and waiting traffic:
\begin{equation}
 \begin{aligned}
 \omega_{\ell cr}&=
 \left[\min(z^+_{\ell r},\zeta_{\ell,c+1})
             -\max(z^-_{\ell r},\zeta_{\ell c})\right]_+,\\
 M_{\ell cm}&=\sum_{r:m_r=m}\omega_{\ell cr},\qquad
 n_\ell=\sum_{c,m}M_{\ell cm}=\sum_r q_{\ell r}.
 \end{aligned}
 \label{eq:compact-grid}
\end{equation}
\par\end{samepage}

The minimum picks the earlier right endpoint and the maximum the later
left endpoint; their positive difference is the shared mass interval.
All these coordinates and overlaps are in vehicle-equivalents, not metres.
This overlap formula is the present representation, not the spatial
CTM flux formula. The grid and ledger describe the same mass, counted
once. A bin containing both left and through mass cannot say which
comes first; the lane ledger supplies that exact sequence for service
and bay access. After transfers, rebuild the grid from the remaining
and admitted ledger portions. Changing mass-bin width changes the representation, not road length
or lane storage.
\par


\head{4. How the task tests information and decisions}
The study uses a progressive mechanism-ablation design: first establish
agreement without order or pocket effects, then isolate shared-lane order
and finite storage, and finally combine storage with ordered bay access.
These controls distinguish a need for order from an effect already
explained by lane counts and available space.

For a local example, take ready $F_1$ traffic $[T(0.4),L(0.6)]$, sending
allowance $1$, and receiving allowances $R_{T_1}=0.2$, $R_B=1$.
Only $0.2$ through mass enters: its residual blocks left access. Reverse
the order without changing lane loads or movement totals, and $0.6$
left mass enters $B$ before $0.2$ through enters $T_1$. Independently
ready $0.3$ in $F_2$ enters $T_2$ in either case. Total entry across the
opening is $0.5$ versus $1.1$. The extra lane remains usable, but does
not erase the inside lane's access effect. Complete-trip consequences
and an aggregate model's ability to infer the obstruction remain questions
to test.

\textbf{Control objective.} Total system time charges each released unit
for every timestep it remains at its source or on a road, until it exits.
Let $o$ index origins and $e$ physical road sections. At the start of
update $k$, $Q_{o,k}$ is mass already waiting at origin $o$ and
$D_{o,k}$ the new release available in that update. Resident road mass is
$n_{e,k}$. The common clearance horizon contains $H$ updates. The objective
is the number of units still in the system, multiplied by their time
charge and summed over updates.
\rev{This reuses the project's source-plus-road population-time quadrature~\cite{m0}, applied to the present section graph:}
\par
\begin{equation}
 J_{\Delta t}
 =\Delta t\sum_{k=0}^{H-1}
 \left[\sum_o(Q_{o,k}+D_{o,k})+\sum_en_{e,k}\right].
 \label{eq:compact-tst}
\end{equation}

Thus $Q+D$ counts source availability and $n$ counts all travelling or
queued road mass. Their units are vehicle-equivalents; multiplying by
$\Delta t$ and summing gives vehicle-equivalent seconds. Admission moves
mass from source to road, so it does not itself remove that mass from
the account. Only sink exit does. Fixed releases are charged before
that update's admission, and all released trips remain included through
a common complete-clearance horizon.

\textbf{Mechanism ablations and robustness.} The scenario codes below
refer to physical tests, not names of comparison models. Order-only and
storage-only controls are separate branches; the combined case tests
their interaction. The final family tests generalization rather than
adding another mechanism.
\begin{center}
\setlength{\tabcolsep}{4pt}
\begin{tabular}{@{}>{\raggedright\arraybackslash}p{0.06\linewidth}>{\raggedright\arraybackslash}p{0.43\linewidth}>{\raggedright\arraybackslash}p{0.43\linewidth}@{}}
\toprule
Case & Mechanism included & What the comparison should establish\\
\midrule
V0 & Movement-exclusive paths; no short bay & Neutral agreement in propagation and service.\\
V1 & Shared-lane order without a finite bay & Whether sequence changes usable service.\\
V2 & Finite bay with order obstruction absent/nonbinding & Whether lane storage explains the effect.\\
V3 & Same bay geometry with obstructing versus accessible order & Whether order adds to the storage explanation.\\
V4 & Reserved variations of the preceding families & Whether conclusions survive unseen conditions.\\
\bottomrule
\end{tabular}
\end{center}
V2 and V3 have the same representative geometry and permissions.
V2 uses arrival patterns for which the bay can be reached whenever room
permits; V3 changes order to make that access consequential. V0 needs
matched service and receiving. Through/right mixing alone need not create
an order effect when both movements have identical usable rates.

``Reserved'' in V4 means scenarios excluded from model calibration and
method tuning. The reserved variables are demand intensity, turn shares,
bay storage/length, arrival-order patterns and downstream receiving,
including asymmetric approaches and combinations not used in the common
symmetric cases. For example, higher left demand with a shorter bay and
reduced downstream room is a distinct reserved combination. Numerical
values and scenario-bank membership must be specified before that test;
these are the variables to reserve, not a claim that a bank has been run.

\textbf{What the simpler models can explain.} The same ablations are
predicted by aggregate link models, structured movement-queue models,
lane/storage aggregate models and the selective ordered model. The
question is whether a simpler traffic description already explains the
response, or whether keeping the realized sequence improves the control
prediction. Aggregate models may derive their states from the original
traffic observations and histories; successful aggregate inference is a
valid alternative explanation, not information that must be withheld.

A movement-queue model, such as Wei et al.'s link queue approach~\cite{wei},
provides a queue-based explanation of waiting and service. Lane-aware
allocation~\cite{gong} and finite-storage models~\cite{liu} provide
stronger explanations when lane use or bay room matters. They are relevant
precisely because V2 should be explainable without exact local order;
V3 asks what remains after those effects are accounted for. The particular
finite-bay composition is an adaptation to be evaluated, not a result
inherited merely by combining citations. The local ordered-mass model
uses continuous lane assignment too, but adds realized sequence and
access to that aggregate account.
\rev{These citations identify comparison-model directions; they do not supply the ordered-model formulas above.}

First compare predicted responses to the same feasible service changes,
then assess whether optimization using those responses improves complete
population time. A richer reference simulation uses the same released
routes and one common mapping from effective service to executable
signals. Its population time includes waiting to enter the network plus
travel duration~\cite{sumo}. All demand is included. This distinguishes
a useful control direction from an internally differentiable model or
an objective that falls only inside its own approximation.
\par


\head{5. How a service change propagates to total time}
Let $\alpha$ be a small amount of service transferred along a feasible
direction $d$, so $g(\alpha)=g+\alpha d$ and $\sum_p d_p=0$.
For example, increasing the left fraction and decreasing the through
fraction by the same amount defines one such control knob. Derivatives
below are written $\mathrm d(\cdot)/\mathrm d\alpha$.

The complete causal path is:
\begin{center}
\begin{tabular}{c}
green allocation $\longrightarrow$ head rates and receiving shares\\
$\downarrow$\\
service rates and available time $\longrightarrow$ discharged mass\\
$\downarrow$\\
remaining ledger masses $\longrightarrow$ lane grid and loads
$\longrightarrow$ actual admissions\\
$\downarrow$\\
entry/readiness times and histories $\longrightarrow$ future sending,
receiving and access\\
$\downarrow$\\
later discharge and sink exits $\longrightarrow$ complete total system time
\end{tabular}
\end{center}
The following paragraphs follow these arrows. Current service changes
first alter discharge; the resulting masses and timing become the state
used by later updates. That later state can change which packet is at a
lane head, which bay is reachable and how much traffic can enter.

Collect source masses, road counts and continuous ledger mass/timing data
in $X_k$. The grid is reconstructed inside the step map $F_k$, assembled
from \eqref{eq:compact-sending}--\eqref{eq:compact-grid}.
Let $g_k$ be the allocation applied in update $k$.
\rev{The first line below defines the composed timestep map; the second follows from it by the chain rule.}
\begin{equation}
 \begin{aligned}
 X_{k+1}&=F_k(X_k,g_k),\\
 \frac{\mathrm dX_{k+1}}{\mathrm d\alpha}
 &=\frac{\partial F_k}{\partial X_k}
       \frac{\mathrm dX_k}{\mathrm d\alpha}
  +\frac{\partial F_k}{\partial g_k}
       \frac{\mathrm dg_k}{\mathrm d\alpha}.
 \end{aligned}
 \label{eq:compact-chain}
\end{equation}
The first term carries earlier queue/history changes into this update;
the second adds the direct effect of current service. During control,
$\mathrm dg_k/\mathrm d\alpha=d$. During fixed clearance it is zero,
but the altered traffic state continues to affect departures.
Common preparation has zero candidate response. The module derivatives
below explain how the two terms of this same chain are formed.

\textbf{Current discharge rates: where green acts directly.}
Differentiating the signal-limited branch of \eqref{eq:compact-action}
gives $\mathrm d\mu_{\ell m}/\mathrm d\alpha=s_{\ell m}d_{p(m)}$;
the fixed physical-cap branch gives zero. Green also changes the leg
weights in \eqref{eq:compact-leg-sharing}. Earlier traffic changes
downstream receiving, with
$\mathrm d\rho_j/\mathrm d\alpha=(\mathrm dR_{j,k}/\mathrm d\alpha)/\Delta t$.
This is the first arrow in the path: control and available room jointly
determine current service rates.

\par\begin{samepage}
Recall that $r_{ij}$ is the rate granted to inbound leg $i$ toward output
$j$, $C_{ij}$ its cap, $w_{ij}$ its green weight and $W_j$ the total usable
output rate. The multiplier $\theta_j$ in \eqref{eq:compact-leg-sharing}
scales uncapped legs' weights. Let ${\cal F}$ contain those uncapped legs
and ${\cal C}$ the capped legs.
\rev{Direct differentiation of \eqref{eq:compact-leg-sharing}, with unchanged cap categories, gives its total-rate balance and categorical responses:}
\begin{equation}
 \begin{aligned}
 \frac{\mathrm dW_j}{\mathrm d\alpha}
 &=\sum_i\frac{\mathrm dr_{ij}}{\mathrm d\alpha},\\
 \frac{\mathrm dr_{ij}}{\mathrm d\alpha}
 &=\begin{cases}
 \displaystyle\frac{\mathrm dC_{ij}}{\mathrm d\alpha},&i\in{\cal C},\\
 \displaystyle w_{ij}\frac{\mathrm d\theta_j}{\mathrm d\alpha}
       +\theta_j\frac{\mathrm dw_{ij}}{\mathrm d\alpha},&i\in{\cal F}.
 \end{cases}
 \end{aligned}
 \label{eq:compact-sharing-derivative}
\end{equation}
\par\end{samepage}

A capped leg follows changes in its cap. An uncapped leg responds both
to the common multiplier and to its own weight. Substituting these cases
into \rev{the balance in \eqref{eq:compact-sharing-derivative} and solving algebraically for the multiplier response} gives
\begin{equation}
 \frac{\mathrm d\theta_j}{\mathrm d\alpha}
 =\frac{\displaystyle
 \frac{\mathrm dW_j}{\mathrm d\alpha}
 -\sum_{i\in{\cal C}}\frac{\mathrm dC_{ij}}{\mathrm d\alpha}
 -\theta_j\sum_{i\in{\cal F}}\frac{\mathrm dw_{ij}}{\mathrm d\alpha}}
 {\displaystyle\sum_{i\in{\cal F}}w_{ij}}.
 \label{eq:compact-sharing-level-derivative}
\end{equation}
Start with the change in total service, subtract what capped legs absorb,
then subtract the rate change already caused by uncapped legs' changing
weights. The remaining response is shared through their total weight.
Compute this multiplier response first, then the individual responses in
\eqref{eq:compact-sharing-derivative}. The same balance differentiates
equal lane sharing in \eqref{eq:compact-lane-sharing}.
If all competitors are capped, their rates follow their caps directly;
the multiplier need not be differentiated by division.

\textbf{Available service time: how order changes the timestep amount.}
Rates alone do not determine discharge: a movement behind a head block
can use only the time left after that block clears.
\rev{Applying the product rule to the amount sum \eqref{eq:compact-service-amount} derives}
\begin{equation}
 \frac{\mathrm dy}{\mathrm d\alpha}
 =\sum_a\left(\delta_a\frac{\mathrm dr_a}{\mathrm d\alpha}
              +r_a\frac{\mathrm d\delta_a}{\mathrm d\alpha}\right).
 \label{eq:compact-service-derivative}
\end{equation}
The first term changes service rate over the existing duration; the
second changes available time at the existing rate. This is the next
arrow in the path: rate and duration responses become actual transferred
mass, which updates the traffic state.

Equation \eqref{eq:compact-head-time} selects the earliest stopping
condition. Consider specifically its \emph{head-clearance branch}, where
the controlling lane has remaining head mass $q$, positive rate $r$ and
$\delta=q/r$.
\rev{The derivative below follows from the selected mass-balance identity $r\delta=q$ in that branch of \eqref{eq:compact-head-time}:}
\begin{equation}
 \frac{\mathrm d\delta}{\mathrm d\alpha}
 =\frac1r\frac{\mathrm dq}{\mathrm d\alpha}
      -\frac q{r^2}\frac{\mathrm dr}{\mathrm d\alpha}.
 \label{eq:compact-head-time-derivative}
\end{equation}
More head mass needs more time; faster service clears it earlier.
The minima no longer appear because this formula assumes the same head
clearance remains the earliest condition. It supplies the duration term
in \eqref{eq:compact-service-derivative} for that branch only.
If sending allowance is exhausted first, replace $q$ by the remaining
allowance. If the timestep ends first, duration is its remaining time:
its derivative is the negative sum of earlier interval-duration changes.
These other branches are written explicitly in the companion.

For V1's shared left/through lane, a left-turn block can prevent through
traffic from using its offered service. Transferring green from through
to left lowers the through rate, but can clear the blocker earlier and
leave more time for through discharge. The time gain can outweigh the
rate loss. This is a concrete use of the two channels in
\eqref{eq:compact-service-derivative}: movement order can change the
response to the same control transfer. A fully cleared head of fixed
mass need not itself discharge more; its followers benefit from earlier
access. The task tests whether retaining that coupling helps predict
control responses beyond aggregate explanations. Its complete-TST value
is not established by this local mechanism; the optional calculation is
in the companion's shared-lane diagnostic.

\textbf{Remaining mass and lane loads: carrying the response forward.}
The path now moves from current discharge to the next traffic state.
The paired update \eqref{eq:compact-update} and continuous ledger splits
retain what did not leave: if packet $r$ serves $y_r$ from resident mass
$q_r$, its remainder responds as $\mathrm d(q_r-y_r)/\mathrm d\alpha$.
Route IDs, route positions and realized order stay structural; continuous
masses determine the shifted cumulative-rank boundaries.
\rev{Differentiating the project representation \eqref{eq:compact-grid}, with fixed bin edges, derives the overlap and lane-mass responses below:}
\begin{equation}
 \frac{\mathrm dM_{\ell cm}}{\mathrm d\alpha}
 =\sum_{r:m_r=m}\frac{\mathrm d\omega_{\ell cr}}{\mathrm d\alpha},
 \qquad
 \frac{\mathrm dn_\ell}{\mathrm d\alpha}
 =\sum_{c,m}\frac{\mathrm dM_{\ell cm}}{\mathrm d\alpha}.
 \label{eq:compact-grid-derivative}
\end{equation}
Fixed bin edges do not move; packet boundaries do. These derivatives
turn changed ledger masses into the lane loads used in later choices.
The grid carries the response of the same physical traffic, not a second
population. No rounding or differentiation through route labels is needed.

\textbf{Actual lane admission: changed loads and room change later entry.}
Recall $V,x_\ell,h_\ell,u_\ell,\lambda$ from
\eqref{eq:compact-allocation}: total admitted mass, its lane additions,
loads before this allocation, admission bounds and the common filling
level. Let ${\cal A}$ contain positive uncapped additions
($0<x_\ell<u_\ell$), and ${\cal U}$ lanes on binding upper bounds.
``Uncapped'' means the bound is not currently binding, not unlimited room.
\rev{Differentiating the total-mass constraint and capped lane rule in \eqref{eq:compact-allocation} derives}
\begin{equation}
 \begin{aligned}
 \frac{\mathrm dV}{\mathrm d\alpha}
 &=\sum_\ell\frac{\mathrm dx_\ell}{\mathrm d\alpha},\\
 \frac{\mathrm dx_\ell}{\mathrm d\alpha}
 &=\begin{cases}
 \displaystyle\frac{\mathrm d\lambda}{\mathrm d\alpha}
             -\frac{\mathrm dh_\ell}{\mathrm d\alpha},&\ell\in{\cal A},\\
 \displaystyle\frac{\mathrm du_\ell}{\mathrm d\alpha},&\ell\in{\cal U},\\
 0,&\text{no addition on the unchanged lower branch}.
 \end{cases}
 \end{aligned}
 \label{eq:compact-allocation-derivative}
\end{equation}
An uncapped addition is the gap between the filling level and starting
load; a capped addition follows its bound; a lane remaining outside the
allocation contributes zero. Substituting these cases into the mass
balance in \rev{\eqref{eq:compact-allocation-derivative} and solving algebraically for the common filling-level sensitivity} gives
\begin{equation}
 \frac{\mathrm d\lambda}{\mathrm d\alpha}
 =\frac1{|{\cal A}|}\left(
 \frac{\mathrm dV}{\mathrm d\alpha}
 -\sum_{\ell\in{\cal U}}\frac{\mathrm du_\ell}{\mathrm d\alpha}
 +\sum_{\ell\in{\cal A}}\frac{\mathrm dh_\ell}{\mathrm d\alpha}\right).
 \label{eq:compact-allocation-level-derivative}
\end{equation}
The first term is the change in total \emph{admitted} mass. Capped lanes'
bound changes absorb part of that total, so they are subtracted.
Higher starting loads reduce uncapped additions at a fixed level, so
their changes are added to find the compensating level response.
Divide by the number of positive uncapped lanes because each responds
equally to a level rise. Compute the level response, then the lane
responses in \eqref{eq:compact-allocation-derivative}.
If no lane is free, use the capped/zero branches without this division.

This is the admission arrow in the larger loop: earlier service changes
lane loads, histories and accessible room, which change where later mass
can enter. The additions update loads and remaining sending/receiving
before the next packet is examined. Retained mass responds as
$\mathrm d(q-V)/\mathrm d\alpha$. Admission has no direct green input
at a fixed starting state; its control response arrives through this
earlier traffic and its effects on the allocation inputs.

\textbf{Entry timing and propagation: when admitted mass becomes usable.}
Changed entry amounts also change the admitted stream's timing.
Recall from \eqref{eq:compact-entry} that $b$ is total section inflow
and $z$ the mass admitted before a particular packet. With fixed timestep
times and $b>0$,
\rev{the quotient-rule derivative of that front-entry endpoint gives}
\begin{equation}
 \frac{\mathrm dt^{\rm in}_{\rm front}}{\mathrm d\alpha}
 =\Delta t\left(\frac1b\frac{\mathrm dz}{\mathrm d\alpha}
       -\frac z{b^2}\frac{\mathrm db}{\mathrm d\alpha}\right).
 \label{eq:compact-entry-derivative}
\end{equation}
The tail follows the same rule with $z+q^{\rm in}$.
This carries the response from lane entry to future availability:
shifting entry/readiness times can change how much of a packet has
travelled far enough by a later update. Fixed free-flow time shifts both
endpoints without changing their derivatives.
The delayed queries in \eqref{eq:compact-sending} and
\eqref{eq:compact-receiving} differentiate stored cumulative counts at
fixed query times. Travel-limited sending follows eligible-mass changes;
storage-limited receiving follows delayed-exit and prior-entry changes.
Those altered limits feed back into later admissions and service.

\textbf{Sink exits and TST: accumulating the control consequence.}
The final arrow translates all those repeated updates into the task
objective. Let $D_{\leq k}$ be all released mass available through update
$k$ and $Z_k$ cumulative sink exits before it.
The \rev{conservation account implied by \eqref{eq:compact-update}} makes
source-plus-road population in \eqref{eq:compact-tst} equal to
$D_{\leq k}-Z_k$. With releases, $\Delta t$ and the common horizon $H$
fixed, \rev{direct differentiation of this conservation-based rewrite of \eqref{eq:compact-tst}} gives
\begin{equation}
 J_{\Delta t}=\Delta t\sum_{k=0}^{H-1}(D_{\leq k}-Z_k),\qquad
 \frac{\mathrm dJ_{\Delta t}}{\mathrm d\alpha}
 =-\Delta t\sum_{k=0}^{H-1}\frac{\mathrm dZ_k}{\mathrm d\alpha}.
 \label{eq:compact-objective-derivative}
\end{equation}
Earlier additional exits reduce population time. Unlike the local
discharge illustration, this accounts for source waiting, moving and
queued traffic through the common complete-clearance horizon, including
the state response after control ends.
The derivatives apply within unchanged head and binding-constraint
patterns; switches can create kinks. Numerical derivative accuracy and
richer-reference control improvement remain distinct evidence from this
derivation. The \href{urban_methodology_specification_formula_provenance.pdf}{mathematical companion}
gives the branch details and optional hand calculation.

\makeatletter
\renewcommand{\section}{\@ifstar{\head}{\head}}
\makeatother
\begin{thebibliography}{10}
\setlength{\itemsep}{0pt}\setlength{\parskip}{0pt}
\bibitem{m0}
{I. Yperman, \emph{The Link Transmission Model for Dynamic
Network Loading}, PhD thesis, KU Leuven, 2007.
\href{https://www.researchgate.net/publication/28360292_The_Link_Transmission_Model_for_dynamic_network_loading}{Author-uploaded thesis}.
The discrete transcription and population-time account here follow the
\href{../../../../docs/milestones/milestone_0_ltm_baseline.md}{repository LTM specification}.}
\bibitem{m1}
{Continuous ordered route-mass and affine timing construction:
\href{../../../../docs/milestones/milestone_1_ordered_route_mass_ledger.md}
{technical representation specification}. The urban allocation/access
composition is defined in the present model.}
\bibitem{gong}
{X. Gong, M.C.J. Bliemer and M.P.H. Raadsen,
\emph{Extended macroscopic node model for multilane traffic},
Transportation Research B 195, 103208, 2025.
\href{https://doi.org/10.1016/j.trb.2025.103208}{doi:10.1016/j.trb.2025.103208}.}
\bibitem{han}
{K. Han et al., \emph{On the continuum approximation of the
on-and-off signal control on dynamic traffic networks}, Transportation
Research B 61, 73--97, 2014.
\href{https://arxiv.org/abs/1309.3461}{Author manuscript};
doi:10.1016/j.trb.2014.01.001.}
\bibitem{node}
{C.M.J. Tamp\`ere et al., \emph{A generic class of first order
node models for dynamic macroscopic simulation of traffic flows},
Transportation Research B 45, 289--309, 2011.
\href{https://doi.org/10.1016/j.trb.2010.06.004}{doi:10.1016/j.trb.2010.06.004}.}
\bibitem{tracking}
{F. de Souza, O. Verbas, J. Auld and C.M.J. Tamp\`ere,
\emph{A mesoscopic link-transmission-model able to track individual vehicles},
Simulation Modelling Practice and Theory, 103088, 2025.
\href{https://www.osti.gov/servlets/purl/2569501}{Open manuscript};
doi:10.1016/j.simpat.2025.103088.}
\bibitem{ctm}
{C.F. Daganzo, \emph{The cell transmission model: A dynamic
representation of highway traffic consistent with the hydrodynamic theory},
Transportation Research B 28, 269--287, 1994.
\href{https://doi.org/10.1016/0191-2615(94)90002-7}{doi:10.1016/0191-2615(94)90002-7}.}
\bibitem{wei}
{Wei et al., link queue model,
\href{https://arxiv.org/html/2310.12249v2}{arXiv:2310.12249v2};
doi:10.1016/j.trc.2024.104937.}
\bibitem{liu}
{Liu and Chang, lane-group storage/blockage model,
Transportation Research C 19, 130--144, 2011.
\href{https://attap.umd.edu/wp-content/uploads/2022/05/TRC522.pdf}{Author-hosted article};
doi:10.1016/j.trc.2010.04.005.}
\bibitem{sumo}
{SUMO official documentation:
\href{https://sumo.dlr.de/docs/Simulation/Traffic_Lights.html}{traffic-light programs}
and \href{https://sumo.dlr.de/docs/Simulation/Output/TripInfo.html}{TripInfo}.}
\end{thebibliography}
\end{document}
